% =====================================================================
\chapter{Trabalhos Relacionados}\label{chap:relatedwork}
% =====================================================================

%\input{chapters/TRABALHOS_RELACIONADOS/trabalhos_decentralizado.tex/}


%\section{Related Work}
Related works found in the literature can be divided into: (a) works that implement Cholesky-based inversion in FPGA or ASIC in the context of MIMO technology and (b) works that implement decentralized equalization in FPGA.

\section{Cholesky-Based Inversion in MIMO Hardware}
The TL-BD-INSA detector of ~\cite{wang2021tlbdinsa}  implements Cholesky-based inversion
of Hermitian block-diagonal sub-matrices on a Virtex-7 XC7VX980T FPGA
and a TSMC 45~nm ASIC. The FPGA implementation targets a $128 \times 8$ MIMO configuration with 64-QAM. The
inversion module uses a triangular systolic array for the Cholesky
factorization, followed by parallel forward and backward substitution
units, one per user equipment. This architecture is the closest available
hardware reference for the $M_c \times M_c$ inversion required by
BCD-MMSE, differing only in that TL-BD-INSA operates under AWGN and
inverts block-diagonal sub-matrices of the Gram matrix, whereas BCD-MMSE
inverts the full local noise-weighted matrix under colored noise.

Table~\ref{tab:hw_cholesky} reports the implementation metrics for the
complete TL-BD-INSA detector pipeline, which includes matched filtering,
Gram matrix computation, the Cholesky inversion module, and symbol
estimation~\cite{wang2021tlbdinsa}. These figures therefore represent
an upper bound on the resources attributable to the inversion step alone;
the paper does not provide a per-module area or resource breakdown.

\begin{table}[h]
\centering
\caption{TL-BD-INSA full detector pipeline hardware results
($128 \times 8$, 64-QAM). FPGA: Virtex-7 XC7VX980T.
ASIC: TSMC 45~nm(~\cite{wang2021tlbdinsa}).}
\label{tab:hw_cholesky}
\begin{tabular}{lcc}
\hline
\textbf{Metric} & \textbf{FPGA} & \textbf{ASIC} \\
\hline
Gate count / LUT slices     & 13,640 LUT  & 572 kGE       \\
Flip-flop slices            & 19,460 FF   & ---           \\
DSP48 units                 & 277         & ---           \\
Clock frequency [MHz]       & 310         & 1,000         \\
Throughput                  & 57.3 Mbps   & 185 Mbps      \\
Power [W]                   & ---         & 0.144         \\
Energy efficiency           & 1,731 bps/slice & 1.28 Gbps/W \\
\hline
\end{tabular}
\end{table}

\section{Decentralized Detection on FPGA}
The FD-LAMA detector of~\cite{li2019fdlama}, implemented on a
Virtex-7 XC7VX690T using High-Level Synthesis (HLS) for a $128 \times 8$
system with QPSK and i.i.d. Rayleigh fading, achieves 22.2~Mb/s
throughput at 429~MHz with 11,673 LUTs and 213 DSP48 blocks for
$C = 1$ cluster. Resource utilization scales roughly linearly with $C$,
while throughput degrades only slightly, reaching 17.8~Mb/s at
$C = 4$~\cite{li2019fdlama}. This work assumes AWGN and does not
include a matrix inversion module: FD-LAMA avoids explicit inversion
entirely, operating only on local Gram matrix-vector products within
each cluster.

The decentralized Newton-method detector of~\cite{kulkarni2021newton},
implemented on a Virtex-7 XC7VX690T for a $128 \times 8$ system with
$C = 4$ clusters in a ring topology, is architecturally equivalent to
the BCD-MMSE daisy-chain in its inter-DU communication pattern. It
achieves 253~Mbps throughput at 340~MHz, consuming 100,733 LUTs,
1,280 DSP48 blocks, and 7.248~W on-chip power for 16-QAM with
$T = 3$ iterations~\cite{kulkarni2021newton}. The ring topology
maintains constant interconnect bandwidth as the MIMO configuration
scales, a property shared by the BCD-MMSE daisy-chain. As with all
other decentralized hardware works, this implementation assumes AWGN
and provides no inversion metrics specific to the colored-noise regime.
\section{Gap Synthesis}

Algorithmic papers also explicitly identify FPGA or ASIC
implementation as future work: \cite{jeon2019} states that decentralized
solutions are expected to be deployed on FPGA or ASIC clusters for higher
computing efficiency and lower transfer latency, and~\cite{li2019cd} notes
that multi-FPGA and multi-ASIC designs are left for future work to reduce
power relative to GPU baselines.

None of the surveyed hardware works simultaneously addresses the three
conditions that define the gap motivating this work: (1)~colored noise
with a non-diagonal local covariance $\hat{\mathbf{R}}_{cc}$;
(2)~decentralized LMMSE equalization under the BCD-MMSE framework; and
(3)~FPGA implementation with measured resource metrics. The algorithmic
solution exists in~\cite{zhao2025jsac}, which proves optimality within the decentralized constraint and
near-centralized BER performance for BCD-MMSE but provides no LUT count,
power, or latency figures for any target platform. This work fills that
gap by providing the first hardware analysis of the $M_c \times M_c$
Hermitian Cholesky inversion at the core of the BCD-MMSE update.
